\chapter{Metodologia}
\label{cap_metodologia}

Este capítulo formaliza os fundamentos metodológicos que regem o SIGEA-GTT, articulando as diretrizes epidemiológicas canônicas do \textit{Global Trigger Tool} (IHI-GTT), o modelo pedagógico de simulação clínica de prontuários fictícios em ambiente isolado, a fundamentação teórica das seis Ferramentas da Qualidade hospitalar e as diretrizes de Engenharia de Software aplicadas à construção do sistema.

\section{A Metodologia IHI-GTT Canônica}
\label{sec_ihi_gtt_metodologia}

O método \textit{Global Trigger Tool}, concebido pelo \textit{Institute for Healthcare Improvement} \cite{griffin2009ihi}, fundamenta-se na premissa de que a auditoria retrospectiva direcionada por marcadores clínicos objetivos supera as limitações de tempo e viabilidade da leitura linear exaustiva de prontuários médicos.

\subsection{Protocolo de Amostragem e Regra dos 20 Minutos}
O protocolo padrão do IHI preconiza a seleção aleatória de **20 prontuários hospitalares por mês**, auditados em ciclos quinzenais de 10 prontuários. Os critérios de elegibilidade amostral exigem que o paciente tenha permanecido internado por no mínimo 24 horas e que a alta, transferência ou óbito já tenha ocorrido (caráter retrospectivo estrito).

A auditoria de cada caso clínico é estritamente delimitada a um **tempo máximo de 20 minutos por prontuário**. Esse limite cronometrado é essencial para garantir a sustentabilidade operacional do método nas instituições de saúde, impedindo que o auditor se perca em minúcias irrelevantes e focando sua atenção na varredura rápida de seções estratégicas: sumário de alta, folhas de evolução clínica e de enfermagem, prescrições farmacológicas e relatórios de exames laboratoriais \cite{resende2019gtt}.

\subsection{Módulos Assistenciais e os 53 Gatilhos Clínicos Oficiais}
O catálogo oficial do IHI estrutura o rastreamento em seis módulos assistenciais interdependentes, totalizando 53 gatilhos canônicos (\autoref{tab_modulos_gtt}):

\begin{table}[htb]
\centering
\caption{Módulos Assistenciais e Distribuição dos 53 Gatilhos Canônicos IHI-GTT.}
\label{tab_modulos_gtt}
\begin{tabularx}{\textwidth}{l c X}
\toprule
\textbf{Módulo Assistencial} & \textbf{Código} & \textbf{Descrição e Foco Clínico} \\
\midrule
Cuidados Gerais & C & 17 gatilhos voltados à assistência de enfermagem e rotina hospitalar (código azul, queda de leito, lesão por pressão, readmissão em 30 dias). \\
Medicamentos & M & 13 gatilhos focados em antídotos específicos e reações adversas graves (vitamina K, naloxona, flumazenil, hipoglicemia severa). \\
Cirúrgico & S & 11 gatilhos para complicações perioperatórias (retorno não programado ao bloco cirúrgico, hemorragia com politransfusão, lesão incidental de órgãos). \\
Terapia Intensiva (UTI) & I & 4 gatilhos para suporte avançado de vida (readmissão em UTI em 48h, reintubação não planejada, pneumonia associada à ventilação mecânica). \\
Perinatal & P & 5 gatilhos materno-infantis (Apgar menor que 6 no 5º minuto, trauma de parto, histerectomia de emergência pós-parto). \\
Urgência e Emergência & E & 3 gatilhos de acolhimento e triagem (retorno à emergência em menos de 48h com hospitalização, permanência em observação superior a 24h). \\
\bottomrule
\end{tabularx}
\fonte{Adaptado de \citeonline{griffin2009ihi}.}
\end{table}

\subsection{Classificação de Gravidade de Dano (NCC MERP)}
A identificação de um gatilho clínico não significa necessariamente a presença de um evento adverso. Uma vez detectado o gatilho, o revisor deve analisar criticamente as circunstâncias clínicas e aplicar o algoritmo padronizado pelo \textit{National Coordinating Council for Medication Error Reporting and Prevention} (NCC MERP) \cite{nccmerp2001}, transcrito na \autoref{tab_ncc_merp}.

\begin{table}[htb]
\centering
\caption{Índice de Categorização de Gravidade e Dano NCC MERP.}
\label{tab_ncc_merp}
\begin{tabularx}{\textwidth}{c c X}
\toprule
\textbf{Dano Real} & \textbf{Categoria} & \textbf{Definição Operacional Padronizada} \\
\midrule
\multirow{4}{*}{\textbf{Sem Dano}} 
& A & Circunstâncias ou eventos com potencial de causar erro ou dano. \\
& B & O erro ocorreu, mas não atingiu o paciente (\textit{near-miss}). \\
& C & O erro atingiu o paciente, mas não causou dano assistencial. \\
& D & O erro atingiu o paciente e exigiu monitoramento para assegurar ausência de dano. \\
\midrule
\multirow{5}{*}{\textbf{Com Dano (EA)}} 
& \textbf{E} & Dano temporário ao paciente que exigiu intervenção clínica ativa. \\
& \textbf{F} & Dano temporário que causou prolongamento da internação hospitalar. \\
& \textbf{G} & Dano permanente com sequela anatômica, funcional ou psicológica. \\
& \textbf{H} & Dano grave com risco iminente de morte, exigindo intervenção para suporte de vida. \\
& \textbf{I} & Óbito do paciente decorrente do evento adverso assistencial. \\
\bottomrule
\end{tabularx}
\fonte{Adaptado de \citeonline{nccmerp2001}.}
\end{table}

Para a metodologia IHI-GTT, **apenas as Categorias E, F, G, H e I caracterizam Eventos Adversos com dano real**, diferenciando o método das ferramentas meramente focadas em quase-falhas (\textit{near-misses}).

\subsection{Indicadores Epidemiológicos Oficiais do IHI}
O IHI-GTT estabelece três métricas matemáticas universais para acompanhamento longitudinal das taxas de incidentes hospitalares:
\begin{enumerate}
    \item \textbf{Taxa de Eventos Adversos por 1.000 Pacientes-Dia} ($TI_{1000}$): Medida de densidade de incidência ponderada pelo tempo de exposição ao risco hospitalar:
    \begin{equation}
        TI_{1000} = \left( \frac{\sum \text{Total de Eventos Adversos (Categorias E a I)}}{\sum \text{Total de Pacientes-Dia na Amostra}} \right) \times 1.000
        \label{eq_ti1000}
    \end{equation}
    
    \item \textbf{Taxa de Eventos Adversos por 100 Admissões} ($TI_{100}$): Medida de frequência de danos por paciente admitido:
    \begin{equation}
        TI_{100} = \left( \frac{\sum \text{Total de Eventos Adversos (Categorias E a I)}}{\sum \text{Total de Admissões Hospitalares Auditadas}} \right) \times 100
        \label{eq_ti100}
    \end{equation}
    
    \item \textbf{Percentual de Admissões com pelo menos um Evento Adverso} ($P_{EA}$): Indicador da proporção de pacientes vulneráveis afetados por falhas assistenciais:
    \begin{equation}
        P_{EA} = \left( \frac{N_{\text{admissões com pelo menos 1 EA}}}{N_{\text{total de admissões auditadas}}} \right) \times 100\%
        \label{eq_pea}
    \end{equation}
\end{enumerate}

\section{O Modelo de Simulação Pedagógica do SIGEA-GTT}
\label{sec_modelo_simulacao}

O SIGEA-GTT adota o paradigma da **Simulação Clínica Computacional Baseada em Casos Fictícios Estruturados**. Para assegurar plena aderência à Lei Geral de Proteção de Dados (LGPD - Lei nº 13.709/2018) e aos preceitos bioéticos da Resolução CNS nº 466/2012, o sistema foi desenhado sob o princípio de **isolamento total de dados reais**. Não existe qualquer conexão entre a base de dados do SIGEA-GTT e os servidores clínicos do Hospital Universitário da UFS.

Em vez disso, os docentes criam e configuram casos clínicos de alta verossimilhança diretamente no sistema. Os casos simulados são armazenados no PostgreSQL 16 utilizando colunas com tipo de dado semiestruturado `JSONB`, contemplando:
\begin{itemize}
    \item \textbf{Identificação Demográfica Fictícia}: Nome social do paciente, idade, sexo, número do leito, clínica de internação e data de admissão;
    \item \textbf{Histórico e Notas de Admissão}: Anamnese detalhada, comorbidades prévias e hipóteses diagnósticas;
    \item \textbf{Evolução Clínica e Multiprofissional Cronológica}: Registros diários com carimbo fictício de médicos assistentes, enfermeiros, fisioterapeutas e nutricionistas, simulando as informações em que os gatilhos se encontram ocultos;
    \item \textbf{Prescrições Médicas Aprazadas}: Medicamentos prescritos com dosagens, vias de administração, horários e carimbos de checagem da enfermagem;
    \item \textbf{Painel de Exames Laboratoriais e de Imagem}: Resultados de hemogramas, eletrólitos, coagulogramas e exames radiológicos acompanhados dos respectivos valores de referência normativos;
    \item \textbf{Procedimentos Cirúrgicos e Invasivos}: Descrição de acessos venosos centrais, ventilação mecânica, sondagens vesicais e laudos operatórios.
\end{itemize}

O discente acessa a atividade, visualiza o caso em abas clínicas e inicia a resolução sob um **cronômetro regressivo de 20 minutos**, replicando com fidelidade a pressão temporal exigida na prática de auditoria hospitalar.

\section{Fundamentação das Ferramentas da Qualidade}
\label{sec_ferramentas_qualidade}

Diferenciando-se das ferramentas meramente burocráticas de auditoria, o SIGEA-GTT estabelece como requisito pedagógico obrigatório que o estudante, ao identificar um evento adverso no prontuário, aplique ferramentas formais da gestão da qualidade para compreender as causas raízes e desenhar contramedidas eficazes. Seis ferramentas foram integradas à plataforma:

\subsection{Diagrama de Causa e Efeito de Ishikawa (6M)}
Proposto por \citeonline{ishikawa1982}, o Diagrama de Ishikawa (espinha de peixe) organiza graficamente as hipóteses causais que convergem para a ocorrência do problema central (efeito adverso). No SIGEA-GTT, o discente analisa as causas nos seis eixos clássicos da manufatura e serviços adaptados à saúde:
\begin{itemize}
    \item \textbf{Método}: Protocolos clínicos desatualizados, ausência de conferência de identidade à beira do leito ou falhas na passagem de plantão;
    \item \textbf{Mão de Obra}: Fadiga da equipe de enfermagem, déficits de treinamento técnico ou sobrecarga de leitos por profissional;
    \item \textbf{Material}: Insumos hospitalares com defeito de fabricação, medicamentos sem identificação visual adequada (\textit{Look-Alike/Sound-Alike});
    \item \textbf{Máquina}: Bombas de infusão descalibradas, ventiladores mecânicos sem manutenção preventiva ou falhas em monitores multiparamétricos;
    \item \textbf{Meio Ambiente}: Ruído excessivo em UTI, iluminação deficiente no preparo de medicações ou pisos escorregadios causadores de quedas;
    \item \textbf{Medida}: Monitorização hemodinâmica em intervalos inadequados ou atraso na liberação de exames laboratoriais críticos.
\end{itemize}

\subsection{Matriz GUT (Gravidade, Urgência e Tendência)}
A Matriz GUT, adaptada da metodologia de análise de problemas de \citeonline{kepner1965}, quantifica as falhas assistenciais para orientar a priorização das intervenções. O aluno atribui notas de 1 a 5 para três dimensões:
\begin{itemize}
    \item \textbf{Gravidade ($G$)}: Impacto e severidade do dano potencial caso o problema persista (1 = sem gravidade até 5 = dano catastrófico/mortal);
    \item \textbf{Urgência ($U$)}: Prazo disponível para agir antes que o dano se materialize (1 = pode esperar até 5 = ação imediatamente indispensável);
    \item \textbf{Tendência ($T$)}: Padrão evolutivo da falha caso nenhuma ação seja tomada (1 = não se agravará até 5 = piora imediata e progressiva).
\end{itemize}
O sistema computa o score de priorização em tempo real via $S_{GUT} = G \times U \times T$, gerando pontuações entre 1 e 125 pontos para ordenamento dinâmico.

\subsection{Plano de Ação 5W2H}
O método 5W2H, consolidado por \citeonline{campos1992}, converte diagnósticos teóricos em planos de contingência executáveis. O SIGEA-GTT disponibiliza um formulário dinâmico no qual o estudante detalha:
\begin{enumerate}
    \item \textit{What} (O que será feito): Descrição precisa da ação corretiva ou preventiva;
    \item \textit{Why} (Por que será feito): Justificativa baseada no gatilho ou evento identificado;
    \item \textit{Where} (Onde será feito): Local ou setor hospitalar (enfermaria, bloco cirúrgico, farmácia satélite);
    \item \textit{When} (Quando será feito): Cronograma e prazos estipulados;
    \item \textit{Who} (Quem fará): Responsável direto pela execução;
    \item \textit{How} (Como será feito): Metodologia técnica e etapas operacionais;
    \item \textit{How Much} (Quanto custará): Estimativa orçamentária ou de alocação de recursos institucionais.
\end{enumerate}

\subsection{Ciclo de Melhoria Contínua PDCA / PDSA}
Desenvolvido por Walter Shewhart e popularizado por \citeonline{deming1986}, o ciclo PDCA (\textit{Plan, Do, Check, Act}) estrutura a melhoria contínua dos processos de cuidado através de quatro fases integradas:
\begin{itemize}
    \item \textbf{Plan (Planejar)}: Reconhecimento da falha assistencial, coleta de evidências e planejamento da meta;
    \item \textbf{Do (Fazer/Executar)}: Implementação controlada da contramedida desenhada;
    \item \textbf{Check (Checar/Verificar)}: Mensuração dos indicadores de resultado em comparação com a meta prévia;
    \item \textbf{Act (Agir/Padronizar)}: Consolidação do protocolo assistencial institucional ou redefinição da estratégia caso a meta não tenha sido atingida.
\end{itemize}

\subsection{Matriz SWOT (FOFA) e Painel de Brainstorming}
A Matriz SWOT (Forças, Oportunidades, Fraquezas e Ameaças) permite ao discente avaliar o ambiente interno e externo da unidade hospitalar antes de implantar uma mudança de segurança do paciente. O módulo é complementado por um painel digital de Brainstorming, estimulando a livre ideação de estratégias assistenciais sem julgamentos prematuros.

\section{Engenharia de Software e Stack Tecnológica}
\label{sec_stack_tecnologica}

O desenvolvimento do SIGEA-GTT fundamentou-se em princípios sólidos da Engenharia de Software Moderna \cite{pressman2021, sommerville2019, martin2018}, privilegiando desacoplamento em camadas, tipagem estática rigorosa e automação de testes com cobertura de 100\%:
\begin{itemize}
    \item \textbf{Backend (Java 21 LTS / Spring Boot 3.3.x)}: Arquitetura em camadas Controller-Service-Repository com DTOs, Bean Validation, MapStruct para mapeamento estático e transações atômicas JPA/Hibernate.
    \item \textbf{Segurança da Informação}: Spring Security 6 com autenticação \textit{stateless} via JSON Web Token (JWT) assinado com HMAC-SHA256, senhas criptografadas com salt adaptativo BCrypt e controle de autorização baseado em perfis estritos (\texttt{ADMIN}, \texttt{PROFESSOR}, \texttt{STUDENT}). O domínio de e-mail é estritamente restrito a \texttt{@academico.ufs.br}.
    \item \textbf{Frontend (Angular 18+)}: Arquitetura reativa baseada em \textit{Standalone Components}, \textit{Signals} do Angular para gerenciamento de estado sem \textit{overhead} de Zone.js, novos blocos de controle de fluxo (\texttt{@if}, \texttt{@for}) e estilização via TailwindCSS com Angular Material.
    \item \textbf{Banco de Dados (PostgreSQL 16 LTS)}: Esquema relacional estruturado gerido por migrações versionadas via Flyway, utilizando atributos nativos `JSONB` para armazenamento flexível de prontuários simulados e ferramentas da qualidade.
    \item \textbf{Critério Estrito de Aceite (Quality Gate)}: 100\% de cobertura de código comprovada por Jacoco no backend (259 testes automatizados) e Jest no frontend (372 testes unitários).
\end{itemize}
